\documentclass[12pt, a4paper]{article}


% 1. Cấu hình Tiếng Việt và Font Times New Roman 13pt chuẩn
\usepackage[utf8]{inputenc}
\usepackage[T5]{fontenc}
\usepackage[vietnamese]{babel}
\usepackage{newtxtext,newtxmath}


\makeatletter
\renewcommand{\normalsize}{\@setfontsize\normalsize{13pt}{16pt}}
\makeatother
\normalsize


% 2. Gói Toán học & Ký hiệu bổ trợ
\usepackage{amsmath, amssymb, amsfonts, mathrsfs, amsthm}
\usepackage{graphicx}
\usepackage{fontawesome}
\usepackage{booktabs, tabularx, array, multirow}
\usepackage{enumitem}


% 3. Đồ họa TikZ, Bảng biến thiên & Hình không gian
\usepackage{tikz}
\usetikzlibrary{calc, intersections, angles, quotes, patterns, arrows.meta, 3d}
\usepackage{pgfplots}
\pgfplotsset{compat=1.18}
\usepackage{tkz-euclide}
\usepackage{tkz-tab}


\renewcommand{\vec}[1]{\overrightarrow{#1}}


% 4. Hộp sư phạm (Tcolorbox)
\usepackage[many]{tcolorbox}
\tcbuselibrary{skins, breakable}


\newtcolorbox{pedabox}[1][]{
    enhanced,
    colback=blue!5!white,
    colframe=blue!75!black,
    fonttitle=\bfseries\sffamily,
    title=#1,
    boxrule=0.8pt,
    arc=2mm,
    breakable,
    left=3mm, right=3mm, top=2mm, bottom=2mm
}


\newtcolorbox{scaffoldbox}[1][]{
    enhanced,
    colback=orange!5!white,
    colframe=orange!85!black,
    fonttitle=\bfseries\sffamily,
    title=#1,
    boxrule=0.8pt,
    arc=2mm,
    breakable,
    left=3mm, right=3mm, top=2mm, bottom=2mm
}


\newtcolorbox{aibox}[1][]{
    enhanced,
    colback=green!5!white,
    colframe=teal!80!black,
    fonttitle=\bfseries\sffamily,
    title=#1,
    boxrule=0.8pt,
    arc=2mm,
    breakable,
    left=3mm, right=3mm, top=2mm, bottom=2mm
}


% 5. Căn lề chuẩn văn bản giáo án
\usepackage{geometry}
\geometry{
    a4paper,
    left=25mm,
    right=15mm,
    top=20mm,
    bottom=20mm
}


% 6. Header / Footer chuẩn hóa (BẮT BUỘC CHÍNH XÁC NỘI DUNG NÀY)
\usepackage{fancyhdr}
\usepackage{lastpage}
\pagestyle{fancy}
\fancyhf{}


\fancyhead[L]{\small \textit{Kế hoạch bài dạy Toán 11}}
\fancyhead[R]{\textbf{Trang \thepage/\pageref{LastPage}}}
\fancyfoot[L]{\small \textit{Tổ Toán - Tin}}
\fancyfoot[R]{\small \textit{Trường THCS\&THPT Hồng Vân}}
\renewcommand{\headrulewidth}{0.4pt}
\renewcommand{\footrulewidth}{0.4pt}


% 7. Giãn dòng & Giãn đoạn theo chuẩn Nghị định 30/2020/NĐ-CP
\usepackage{setspace}
\setstretch{1.15}
\setlength{\parindent}{1.0cm}
\setlength{\parskip}{5pt}


\begin{document}


\begin{center}
    {\Large \textbf{KẾ HOẠCH BÀI DẠY MÔN TOÁN LỚP 11}}\\[0.4em]
    {\LARGE \textbf{BÀI 6. CẤP SỐ CỘNG}}\\[0.5em]
    \textbf{Môn học:} Toán 11 | \textbf{Thời lượng:} 02 tiết (Tiết 13 và Tiết 14 theo PPCT)
\end{center}


\vspace{0.5em}
\hrule
\vspace{1em}


\section*{I. MỤC TIÊU DẠY HỌC}


\subsection*{1. Kiến thức, Năng lực}


\textbf{a) Năng lực Toán học đặc thù (nguyên văn chuẩn CT GDPT 2018):}
\begin{itemize}[leftmargin=1.5em]
    \item \textbf{Năng lực tư duy và lập luận toán học:}
    \begin{itemize}
        \item Nhận biết được một dãy số là cấp số cộng thông qua định nghĩa: số hạng sau bằng số hạng trước cộng với một hằng số $d$ không đổi (công sai).
        \item Giải thích được công thức xác định số hạng tổng quát của cấp số cộng $u_n = u_1 + (n - 1)d$ thông qua phương pháp quy nạp toán học hoặc liệt kê truy hồi.
        \item Nhận biết và chứng minh được tính chất đặc trưng của cấp số cộng: Mỗi số hạng (trừ số hạng đầu và số hạng cuối) đều là trung bình cộng của hai số hạng kề nó: $u_k = \dfrac{u_{k-1} + u_{k+1}}{2}$ ($k \ge 2$).
        \item Giải thích và suy dẫn được công thức tính tổng của $n$ số hạng đầu tiên của cấp số cộng: $S_n = \dfrac{n(u_1 + u_n)}{2} = \dfrac{n[2u_1 + (n - 1)d]}{2}$.
    \end{itemize}
    \item \textbf{Năng lực giải quyết vấn đề toán học:}
    \begin{itemize}
        \item Xác định được các yếu tố cơ bản của một cấp số cộng: số hạng đầu $u_1$, công sai $d$, số hạng thứ $n$ ($u_n$), số lượng số hạng $n$ và tổng $S_n$.
        \item Giải được hệ phương trình tìm $u_1$ và $d$ khi biết hai điều kiện tuyến tính của các số hạng.
        \item Tính được tổng $S_n$ của một cấp số cộng trong các bài toán đại số và hình học.
    \end{itemize}
    \item \textbf{Năng lực mô hình hóa toán học:}
    \begin{itemize}
        \item Vận dụng kiến thức cấp số cộng giải quyết các bài toán thực tiễn liên môn: tính số lượng ống thép hoặc khúc gỗ xếp chồng hình thang cân; tính lương tăng theo định kì; tính tổng quãng đường rơi tự do tăng dần đều; bài toán số ghế ngồi tăng theo hàng trong nhà hát/sân vận động.
    \end{itemize}
    \item \textbf{Năng lực giao tiếp toán học:}
    \begin{itemize}
        \item Sử dụng chính xác các thuật ngữ: cấp số cộng, công sai, số hạng tổng quát, tổng $n$ số hạng đầu; trình bày lời giải toán học mạch lạc, rõ ràng từng bước thế công thức.
    \end{itemize}
    \item \textbf{Năng lực sử dụng công cụ, phương tiện học toán:}
    \begin{itemize}
        \item Sử dụng máy tính cầm tay (MTCT Casio fx-580VN X) giải hệ phương trình bậc nhất hai ẩn tìm $u_1$ và $d$; tính nhanh tổng $S_n$ bằng phím $\sum$.
        \item Khai thác bảng tính Excel để tính toán tự động các số hạng và tính tổng cấp số cộng với số lượng phần tử lớn.
    \end{itemize}
\end{itemize}


\textbf{b) Năng lực số (theo Thông tư 02/2025/TT-BGDĐT --- Mã 5.3NC1a):}
\begin{itemize}[leftmargin=1.5em]
    \item Học sinh thiết lập công thức tính tổng cấp số cộng trên bảng tính Excel phòng máy 35 máy bằng cả hai phương pháp: dùng hàm \texttt{=SUM(...)} trên cột dữ liệu tự động sinh và dùng trực tiếp công thức đại số \texttt{=N*(2*U1+(N-1)*D)/2}; so sánh kết quả và kiểm chứng tính chính xác của thuật toán số.
\end{itemize}


\textbf{c) Năng lực AI (theo Quyết định 2422/QĐ-BGDĐT --- Mã 11.C2.2):}
\begin{itemize}[leftmargin=1.5em]
    \item Học sinh đề xuất được tính năng AI phân tích quy luật tăng trưởng tuyến tính tương tự cấp số cộng; giải thích được cơ chế mô hình hồi quy tuyến tính (Linear Regression) trong học máy (Machine Learning) tìm đường thẳng $y = ax + b$ khớp với chuỗi dữ liệu rời rạc có bước nhảy xấp xỉ không đổi $d$.
\end{itemize}


\textbf{d) Năng lực chung \& Phẩm chất:}
\begin{itemize}[leftmargin=1.5em]
    \item \textbf{Năng lực chung:} Tự chủ và tự học (chủ động tìm tòi mối liên hệ giữa các số hạng); giao tiếp và hợp tác nhóm khi cùng thực hành thiết kế mô hình STEM.
    \item \textbf{Phẩm chất:} Cẩn thận, chính xác khi biến đổi đại số (đặc biệt không bỏ quên chia $2$ trong công thức $S_n$); trung thực và trách nhiệm trong học tập.
\end{itemize}


\section*{II. THIẾT BỊ DẠY HỌC VÀ HỌC LIỆU}
\begin{itemize}[leftmargin=1.5em]
    \item \textbf{Giáo viên:} Kế hoạch bài dạy chuẩn CV 5512; bài giảng PowerPoint; tệp bảng tính Excel mẫu minh họa công thức $u_n$ và $S_n$; phòng máy vi tính 35 máy; Phiếu học tập số 1 và số 2.
    \item \textbf{Học sinh:} SGK Toán 11; vở ghi bài; 35 máy tính cầm tay Casio fx-580VN X; mô hình que tính / cốc giấy xếp chồng để trực quan hóa cấp số cộng.
\end{itemize}


\section*{III. TIẾN TRÌNH DẠY HỌC (TRỌN VẸN 02 TIẾT)}


\begin{center}
    \framebox[1.0\textwidth]{\parbox{0.96\textwidth}{
        \textbf{CẤU TRÚC PHÂN PHỐI 02 TIẾT DẠY HỌC:}\\[0.3em]
        $\bullet$ \textbf{Tiết 1 (Tiết 13 PPCT):} Định nghĩa cấp số cộng --- Công sai $d$ --- Công thức số hạng tổng quát $u_n = u_1 + (n - 1)d$ --- Tính chất số hạng --- Tích hợp Năng lực AI hồi quy tuyến tính (Mã 11.C2.2).\\[0.3em]
        $\bullet$ \textbf{Tiết 2 (Tiết 14 PPCT):} Tổng $n$ số hạng đầu tiên của cấp số cộng $S_n$ --- Luyện tập tính toán hệ điều kiện --- Năng lực số Excel tính tổng (Mã 5.3NC1a) --- Chủ đề STEM xếp chồng ống thép.
    }}
\end{center}


\vspace{1em}
\hrule
\vspace{0.8em}


{\large \textbf{TIẾT 1 (TIẾT 13 THEO PHÂN PHỐI CHƯƠNG TRÌNH)}}\\[0.3em]
\textbf{Nội dung:} Định nghĩa cấp số cộng --- Công sai $d$ --- Công thức số hạng tổng quát --- Tính chất trung bình cộng --- Năng lực AI (Mã 11.C2.2).


\subsubsection*{1. Hoạt động 1: Mở đầu / Khởi động (05 phút)}
\textbf{Chủ đề:} \textit{Bài toán thang lương tích lũy theo năm công tác.}
\begin{itemize}
    \item \textbf{a) Mục tiêu:} Tạo tình huống có bối cảnh thực tiễn về sự tăng trưởng đều đặn của một đại lượng sau mỗi mốc thời gian, dẫn dắt tự nhiên vào định nghĩa cấp số cộng.
    \item \textbf{b) Nội dung:} Một kĩ sư trẻ bắt đầu đi làm với mức lương khởi điểm năm thứ nhất là $10$ triệu đồng/tháng. Theo hợp đồng, cứ sau mỗi năm làm việc, mức lương hằng tháng của anh được tăng thêm đều đặn $1{,}5$ triệu đồng.
    Hãy viết mức lương hằng tháng của kĩ sư đó trong $5$ năm đầu tiên và tìm quy luật tính mức lương ở năm thứ $n$.
    \item \textbf{c) Sản phẩm:} Học sinh liệt kê dãy số lương (triệu đồng): $10; 11{,}5; 13; 14{,}5; 16$.
    Học sinh phát hiện: Số liền sau bằng số liền trước cộng thêm $1{,}5$. Ở năm thứ $n$, lương được cộng thêm $(n - 1)$ lần mức $1{,}5$ triệu đồng.
    \item \textbf{d) Tổ chức thực hiện:} GV nêu bài toán $\to$ HS thảo luận cặp đôi $2$ phút $\to$ Đại diện $1$ HS trình bày $\to$ GV nhận xét và kết nối vào khái niệm Cấp số cộng.
\end{itemize}


\subsubsection*{2. Hoạt động 2: Hình thành kiến thức mới (25 phút)}


\paragraph{Mục 1: Định nghĩa cấp số cộng & Công sai (08 phút)}
\begin{itemize}
    \item \textbf{Định nghĩa:} Cấp số cộng là một dãy số (hữu hạn hay vô hạn), trong đó kể từ số hạng thứ hai, mỗi số hạng đều bằng số hạng đứng ngay trước nó cộng với một số không đổi $d$.
    Số $d$ được gọi là \textbf{công sai} của cấp số cộng.
    \item \textbf{Hệ thức truy hồi:}
    $$(u_n) \text{ là cấp số cộng} \iff u_{n+1} = u_n + d, \quad \forall n \in \mathbb{N}^*$$
    \item \textbf{Công thức tính công sai:}
    $$d = u_{n+1} - u_n, \quad \forall n \in \mathbb{N}^*$$
    \item \textbf{Trường hợp đặc biệt:}
    \begin{itemize}
        \item Khi $d = 0$: Cấp số cộng có dạng $u_1, u_1, u_1, \dots$ (dãy số không đổi).
        \item Khi $d > 0$: Cấp số cộng là dãy số tăng.
        \item Khi $d < 0$: Cấp số cộng là dãy số giảm.
    \end{itemize}
\end{itemize}


\begin{scaffoldbox}[Giàn giáo sư phạm cho HS TB - Yếu: Mẹo tính đúng công sai $d$]
    \textbf{Lưu ý quy tắc tránh sai dấu:}
    \begin{itemize}
        \item Công sai $d = \text{(Số đứng sau)} - \text{(Số đứng trước liền kề)}$. Cụ thể: $d = u_2 - u_1 = u_3 - u_2$.
        \item \textbf{Sai lầm thường gặp:} Học sinh yếu hay lấy $u_1 - u_2$ dẫn đến sai dấu của công sai $d$, kéo theo sai toàn bộ bài toán.
        \item Ví dụ: Cho dãy số $8, 5, 2, -1, \dots$ Công sai là $d = 5 - 8 = -3$ (dãy giảm, $d < 0$), KHÔNG PHẢI $d = 8 - 5 = 3$.
    \end{itemize}
\end{scaffoldbox}


\paragraph{Mục 2: Số hạng tổng quát của cấp số cộng (10 phút)}
\begin{itemize}
    \item \textbf{Định lí 1:} Nếu cấp số cộng $(u_n)$ có số hạng đầu $u_1$ và công sai $d$ thì số hạng tổng quát $u_n$ được xác định bởi công thức:
    $$u_n = u_1 + (n - 1)d, \quad \forall n \ge 2$$
    \item \textbf{Giải thích quy luật:}
    \begin{itemize}
        \item $u_2 = u_1 + d$
        \item $u_3 = u_2 + d = (u_1 + d) + d = u_1 + 2d$
        \item $u_4 = u_3 + d = (u_1 + 2d) + d = u_1 + 3d$
        \item \dots
        \item $u_n = u_1 + (n - 1)d$.
    \end{itemize}
    \item \textbf{Tính chất các số hạng:} Với cấp số cộng $(u_n)$, mỗi số hạng kể từ số hạng thứ hai đều là trung bình cộng của hai số hạng kề nó:
    $$u_k = \dfrac{u_{k-1} + u_{k+1}}{2} \iff u_{k-1} + u_{k+1} = 2u_k, \quad \forall k \ge 2$$
\end{itemize}


\paragraph{Mục 3: Tích hợp Năng lực AI (Mã 11.C2.2) (07 phút)}


\begin{aibox}[Tích hợp Năng lực AI (QĐ 2422/QĐ-BGDĐT --- Mã 11.C2.2): AI và Hồi quy tuyến tính]
    $\bullet$ \textbf{Mối liên hệ toán học cốt lõi:}\\
    Công thức số hạng tổng quát của cấp số cộng có dạng:
    $$u_n = d \cdot n + (u_1 - d)$$
    Đây chính là hàm số bậc nhất tuyến tính $y = ax + b$ với hệ số góc $a = d$ và tung độ gốc $b = u_1 - d$.\\
    $\bullet$ \textbf{Ứng dụng trong Trí tuệ nhân tạo (Linear Regression):}\\
    Khi hệ thống AI thu thập các dữ liệu thực tế (như giá nhà theo diện tích, chi phí quảng cáo theo doanh số), các điểm dữ liệu thường rải rác xung quanh một xu hướng tăng trưởng tuyến tính. Thuật toán AI tìm cách xác định đường thẳng hồi quy $y = \hat{a}x + \hat{b}$ tối ưu nhất. Độ dốc $\hat{a}$ của đường thẳng này tương đương chính xác với công sai $d$ của cấp số cộng: phản ánh đại lượng tăng thêm bao nhiêu đơn vị sau mỗi bước nhảy!
\end{aibox}


\subsubsection*{3. Hoạt động 3: Luyện tập (10 phút)}
\begin{itemize}
    \item \textbf{Bài 1:} Cho cấp số cộng $(u_n)$ có số hạng đầu $u_1 = 3$ và công sai $d = 4$.
    \begin{itemize}
        \item a) Tìm số hạng thứ $10$ ($u_{10}$) và số hạng thứ $25$ ($u_{25}$).\\
        \textit{Lời giải:}\\
        Áp dụng công thức $u_n = u_1 + (n - 1)d$:\\
        $u_{10} = 3 + (10 - 1) \cdot 4 = 3 + 9 \cdot 4 = 39$.\\
        $u_{25} = 3 + (25 - 1) \cdot 4 = 3 + 24 \cdot 4 = 99$.
        \item b) Số $199$ có phải là một số hạng của cấp số cộng trên không? Nếu có thì là số hạng thứ mấy?\\
        \textit{Lời giải:}\\
        Giả sử $u_n = 199 \iff 3 + (n - 1) \cdot 4 = 199 \iff 4(n - 1) = 196 \iff n - 1 = 49 \iff n = 50$.\\
        Vì $n = 50 \in \mathbb{N}^*$ nên số $199$ là số hạng thứ $50$ của cấp số cộng ($u_{50} = 199$).
    \end{itemize}
    \item \textbf{Bài 2:} Tìm cấp số cộng $(u_n)$ biết: $\begin{cases} u_1 + u_5 = 18 \\ u_3 + u_7 = 26 \end{cases}$.\\
    \textit{Lời giải:}\\
    Biểu diễn các số hạng theo $u_1$ và $d$:
    $$\begin{cases} u_1 + (u_1 + 4d) = 18 \\ (u_1 + 2d) + (u_1 + 6d) = 26 \end{cases} \iff \begin{cases} 2u_1 + 4d = 18 \\ 2u_1 + 8d = 26 \end{cases}$$
    Trừ vế theo vế: $4d = 8 \iff d = 2 \implies 2u_1 = 18 - 4(2) = 10 \iff u_1 = 5$.\\
    Vậy cấp số cộng có $u_1 = 5$ và $d = 2$.
\end{itemize}


\subsubsection*{4. Hoạt động 4: Vận dụng (05 phút)}
\textbf{Bài toán thực tế âm nhạc (Tần số nốt nhạc):}\\
Trong một dàn chuông gió đồng trục, chuông ngắn nhất có chiều dài $20\text{ cm}$. Chiều dài của các chuông tiếp theo tăng dần theo một cấp số cộng với công sai $d = 2{,}5\text{ cm}$. Dàn chuông gồm có $8$ chiếc. Hỏi chiếc chuông dài nhất có chiều dài bằng bao nhiêu?\\
\textit{Lời giải:}\\
Dãy chiều dài của $8$ chiếc chuông lập thành cấp số cộng với $u_1 = 20\text{ cm}$, $d = 2{,}5\text{ cm}$, $n = 8$.\\
Chiều dài của chiếc chuông dài nhất là số hạng thứ $8$:
$$u_8 = u_1 + (8 - 1)d = 20 + 7 \cdot 2{,}5 = 20 + 17{,}5 = 37{,}5\text{ (cm)}$$


\newpage


{\large \textbf{TIẾT 2 (TIẾT 14 THEO PHÂN PHỐI CHƯƠNG TRÌNH)}}\\[0.3em]
\textbf{Nội dung:} Tổng $n$ số hạng đầu tiên của cấp số cộng ($S_n$) --- Kĩ năng giải hệ tìm tham số --- Năng lực số Excel (Mã 5.3NC1a) --- Chủ đề STEM xếp chồng ống thép.


\subsubsection*{1. Hoạt động 1: Mở đầu / Khởi động (05 phút)}
\textbf{Chủ đề:} \textit{Giai thoại thần đồng toán học Carl Friedrich Gauss.}
\begin{itemize}
    \item \textbf{a) Mục tiêu:} Khơi gợi sự tò mò và cảm hứng sáng tạo về phương pháp ghép cặp đối xứng để tính tổng dãy số cách đều.
    \item \textbf{b) Nội dung:} Thầy giáo tiểu học giao cho cậu bé Gauss $10$ tuổi tính tổng từ $1$ đến $100$: $S = 1 + 2 + 3 + \dots + 99 + 100$.
    Chỉ trong vài giây, Gauss đã đưa ra đáp số chính xác là $5050$. Em hãy giải thích cách làm thông minh của Gauss.
    \item \textbf{c) Sản phẩm:} Học sinh ghép cặp: $(1 + 100) = 101, (2 + 99) = 101, \dots, (50 + 51) = 101$. Có tất cả $50$ cặp số có tổng bằng $101 \implies S = 50 \times 101 = 5050$.
    \item \textbf{d) Tổ chức thực hiện:} GV kể chuyện $\to$ HS tái hiện cách ghép cặp $\to$ GV tổng quát hóa thành công thức tính $S_n$.
\end{itemize}


\subsubsection*{2. Hoạt động 2: Hình thành kiến thức mới (20 phút)}


\paragraph{Mục 1: Công thức tính tổng $n$ số hạng đầu tiên $S_n$ (10 phút)}
\begin{itemize}
    \item Cho cấp số cộng $(u_n)$. Đặt $S_n = u_1 + u_2 + u_3 + \dots + u_n$.
    \item Viết xuôi: $S_n = u_1 + (u_1 + d) + \dots + [u_1 + (n - 1)d]$.
    \item Viết ngược: $S_n = u_n + (u_n - d) + \dots + [u_n - (n - 1)d]$.
    \item Cộng từng vế hai đẳng thức trên:
    $$2S_n = (u_1 + u_n) + (u_1 + u_n) + \dots + (u_1 + u_n) \quad (n \text{ số hạng})$$
    $$2S_n = n(u_1 + u_n) \implies S_n = \dfrac{n(u_1 + u_n)}{2}$$
    \item Mặt khác, vì $u_n = u_1 + (n - 1)d$ nên thay vào ta có:
    $$S_n = \dfrac{n[2u_1 + (n - 1)d]}{2}$$
\end{itemize}


\begin{scaffoldbox}[Giàn giáo sư phạm cho HS TB - Yếu: Khắc sâu 2 công thức tính $S_n$]
    \textbf{Nhận biết tình huống sử dụng công thức:}
    \begin{itemize}
        \item \textbf{Công thức 1:} $S_n = \dfrac{n(u_1 + u_n)}{2}$ (Dùng khi ĐÃ BIẾT số hạng đầu $u_1$ và số hạng cuối $u_n$).
        \item \textbf{Công thức 2:} $S_n = \dfrac{n[2u_1 + (n - 1)d]}{2}$ (Dùng khi CHƯA BIẾT số hạng cuối, chỉ biết $u_1$ và công sai $d$).
        \item \textbf{CẢNH BÁO QUAN TRỌNG:} Tuyệt đối KHÔNG ĐƯỢC QUÊN chia cho $2$. Học sinh rất hay quên chia $2$ dẫn đến kết quả bị gấp đôi!
    \end{itemize}
\end{scaffoldbox}


\paragraph{Mục 2: Tích hợp Năng lực số (Mã 5.3NC1a) (10 phút)}


\begin{pedabox}[Năng lực số (Thông tư 02/2025/TT-BGDĐT --- Mã 5.3NC1a): Lập công thức Excel tính $S_n$]
    $\bullet$ \textbf{Mục tiêu thực hành phòng máy 35 máy:}
    \begin{itemize}
        \item Nhập tham số đầu vào tại ô B1 ($u_1 = 5$), ô B2 ($d = 3$), ô B3 ($n = 100$).
        \item \textbf{Cách 1 (Dùng công thức đại số):} Tại ô B4, nhập công thức:
        $$\texttt{=B3*(2*B1+(B3-1)*B2)/2}$$
        Máy tính lập tức trả về kết quả $S_{100} = 15350$.
        \item \textbf{Cách 2 (Dùng hàm SUM lũy tiến):}
        Tại cột C sinh dãy số tự động từ $u_1$ đến $u_{100}$. Tại ô C101 nhập lệnh \texttt{=SUM(C1:C100)}.
        \item Học sinh đối chiếu kết quả của hai cách để khẳng định tính nhất quán của toán học và thuật toán bảng tính số.
    \end{itemize}
\end{pedabox}


\subsubsection*{3. Hoạt động 3: Luyện tập (12 phút)}
\begin{itemize}
    \item \textbf{Bài 1:} Tính tổng của $30$ số hạng đầu tiên của cấp số cộng $(u_n)$ biết:
    \begin{itemize}
        \item a) $u_1 = -2$ và $u_{30} = 85$.\\
        \textit{Lời giải:} Áp dụng công thức 1:
        $$S_{30} = \dfrac{30(u_1 + u_{30})}{2} = 15 \cdot (-2 + 85) = 15 \cdot 83 = 1245$$
        \item b) $u_1 = 4$ và công sai $d = 3$.\\
        \textit{Lời giải:} Áp dụng công thức 2:
        $$S_{30} = \dfrac{30[2(4) + (30 - 1) \cdot 3]}{2} = 15 \cdot [8 + 29 \cdot 3] = 15 \cdot [8 + 87] = 15 \cdot 95 = 1425$$
    \end{itemize}
    \item \textbf{Bài 2:} Một cấp số cộng $(u_n)$ có $u_1 = 3$ và công sai $d = 5$. Biết $S_n = 390$, hãy tìm số lượng số hạng $n$.\\
    \textit{Lời giải:}\\
    Ta có: $S_n = \dfrac{n[2u_1 + (n - 1)d]}{2} = 390 \iff n[2(3) + (n - 1) \cdot 5] = 780$\\
    $$\iff n[6 + 5n - 5] = 780 \iff n(5n + 1) = 780 \iff 5n^2 + n - 780 = 0$$
    Bấm MTCT Casio fx-580VN X (\texttt{MENU} $9 \to 2 \to 2$):
    Phương trình có nghiệm $n = 12$ hoặc $n = -\dfrac{65}{5} = -13$ (loại vì $n \in \mathbb{N}^*$).\\
    Vậy số lượng số hạng là $n = 12$.
\end{itemize}


\subsubsection*{4. Hoạt động 4: Vận dụng (Chủ đề Giáo dục STEM) (08 phút)}
\textbf{Chủ đề STEM:} \textit{Thiết kế và tính toán xếp chồng ống thép công nghiệp.}\\
Tại một công trường xây dựng, người ta xếp các ống thép tròn thành một khối có mặt cắt hình thang cân. Hàng dưới cùng (tầng đáy) có $30$ ống, mỗi hàng liền trên bớt đi $1$ ống so với hàng dưới nó. Hàng trên cùng có $10$ ống.
\begin{itemize}
    \item a) Khối ống thép đó có tất cả bao nhiêu hàng (tầng)?
    \item b) Tính tổng số lượng ống thép có trong toàn bộ khối xếp đó.
\end{itemize}
\textit{Lời giải:}\\
a) Số lượng ống thép ở các hàng từ trên xuống dưới lập thành một cấp số cộng với số hạng đầu $u_1 = 10$ (hàng trên cùng), công sai $d = 1$, số hạng cuối $u_n = 30$ (hàng đáy).\\
Ta có: $u_n = u_1 + (n - 1)d \iff 30 = 10 + (n - 1) \cdot 1 \iff n - 1 = 20 \iff n = 21$.\\
Vậy khối ống thép có tất cả $21$ hàng.\\
b) Tổng số lượng ống thép trong toàn bộ khối xếp là tổng $S_{21}$ của cấp số cộng:
$$S_{21} = \dfrac{n(u_1 + u_n)}{2} = \dfrac{21 \cdot (10 + 30)}{2} = \dfrac{21 \cdot 40}{2} = 21 \cdot 20 = 420\text{ (ống)}$$
*Kết luận:* Toàn bộ khối có $420$ ống thép được xếp trong $21$ hàng.


\newpage


\section*{IV. PHỤ LỤC HỌC LIỆU BỔ TRỢ}


\subsection*{Phụ lục 1: Phiếu học tập số 1 (Tiết 13) --- Định nghĩa \& Số hạng tổng quát CSC}
\noindent\textbf{Họ và tên:} .................................................................... \textbf{Lớp:} 11A....... \textbf{Nhóm:} ..........


\begin{pedabox}[Nhiệm vụ 1: Hoàn thành thông tin công thức Cấp số cộng]
    $\bullet$ Dãy số $(u_n)$ là cấp số cộng khi và chỉ khi: $u_{n+1} = $ .......................................... ($n \ge 1$).\\[0.3em]
    $\bullet$ Công thức tính công sai từ hai số hạng liền kề: $d = $ .............................................\\[0.3em]
    $\bullet$ Công thức số hạng tổng quát: $u_n = $ .................................................................... ($n \ge 2$).\\[0.3em]
    $\bullet$ \textbf{Bài tập nhanh:} Cho cấp số cộng có $u_1 = -3$ và $d = 5$. Tìm $u_7$.\\
    \textit{Bài làm:} ...........................................................................................................................................
\end{pedabox}


\subsection*{Phụ lục 2: Phiếu học tập số 2 (Tiết 14) --- Tổng $n$ số hạng đầu tiên \& Ứng dụng STEM}
\noindent\textbf{Họ và tên:} .................................................................... \textbf{Lớp:} 11A....... \textbf{Nhóm:} ..........


\begin{pedabox}[Nhiệm vụ 2: Điền khuyết công thức tính tổng $S_n$]
    $\bullet$ Công thức 1 (khi biết $u_1$ và $u_n$): $S_n = $ ........................................................................\\[0.3em]
    $\bullet$ Công thức 2 (khi biết $u_1$ và $d$): $S_n = $ ........................................................................\\[0.3em]
    $\bullet$ \textbf{Bài toán thực tế:} Một rạp hát có $20$ hàng ghế. Hàng đầu tiên có $15$ ghế, các hàng sau mỗi hàng tăng thêm $2$ ghế so với hàng trước. Tính tổng số ghế trong rạp hát.\\
    \textit{Bài làm:} ...........................................................................................................................................\\
    .............................................................................................................................................................
\end{pedabox}


\subsection*{Phụ lục 3: Bảng Rubric đánh giá năng lực tích hợp trong Bài 6}


\begin{center}
\begin{tabularx}{\linewidth}{|p{3.2cm}|X|X|X|X|}
\hline
\textbf{Tiêu chí} & \textbf{Mức 1 (Chưa đạt)} & \textbf{Mức 2 (Đạt)} & \textbf{Mức 3 (Khá)} & \textbf{Mức 4 (Tốt)} \\
\hline
\textbf{1. Giải toán Cấp số cộng} & Tính nhầm công sai $d$ do lấy ngược số hạng; quên chia $2$ trong công thức $S_n$. & Tìm được $u_n$ và tính được $S_n$ với các dữ liệu cho trực tiếp. & Giải thành thạo hệ phương trình tìm $u_1, d$; tìm được số lượng số hạng $n$. & Vận dụng xuất sắc mô hình hóa giải các bài toán thực tiễn STEM phức tạp. \\
\hline
\textbf{2. Ứng dụng Excel (Mã 5.3NC1a)} & Chưa biết thiết lập công thức tính tổng trong bảng tính. & Nhập đúng công thức tính $S_n$ vào một ô tính theo tham số cho sẵn. & Sử dụng thành thạo cả hai cách (công thức đại số và hàm \texttt{SUM}) để đối chứng. & Thiết kế bảng tính chuyên nghiệp, tự động hóa toàn bộ việc tính toán các bài toán cấp số cộng. \\
\hline
\textbf{3. Năng lực AI (Mã 11.C2.2)} & Không nhận ra mối liên hệ giữa hàm số bậc nhất và cấp số cộng. & Nhận biết được cấp số cộng là mô hình rời rạc của hàm số tuyến tính $y = ax + b$. & Trình bày được ứng dụng của hồi quy tuyến tính trong học máy AI dựa trên bước nhảy $d$. & Đề xuất được giải pháp ứng dụng AI phát hiện quy luật tuyến tính trong các tập dữ liệu thực nghiệm. \\
\hline
\end{tabularx}
\end{center}


\end{document}